\documentclass[11pt,a4paper]{article}
% preamble.tex — shared setup for human-verification documents
% Compile each document with:  tectonic <File>.tex

\usepackage[T1]{fontenc}
\usepackage[utf8]{inputenc}
\usepackage{lmodern}
\usepackage{amssymb,amsmath,amsthm}
\usepackage{mathtools}
\usepackage[dvipsnames]{xcolor}
\usepackage{listings}
\usepackage{hyperref}
\usepackage[margin=2.5cm]{geometry}
\usepackage{parskip}

% ── Colours ──────────────────────────────────────────────────────────
\definecolor{leanBg}{HTML}{F7F7F0}
\definecolor{haskellBg}{HTML}{F0F4F7}
\definecolor{keyword}{HTML}{7B3294}
\definecolor{comment}{HTML}{5A8C32}
\definecolor{noteFg}{HTML}{B35900}

% ── Lean 4 listings language ─────────────────────────────────────────
% NOTE: Lean listing blocks use ASCII approximations of unicode symbols
% (e.g. -> for arrow, o. for compose).  The literate key substitutes
% these with proper math-mode renderings in the PDF.
\lstdefinelanguage{Lean4}{
  morekeywords={def,theorem,structure,where,namespace,end,import,fun,
                let,have,show,by,exact,intro,cases,rfl,subst,obtain,
                funext,simp,rw,apply,variable,Type,Prop,instance,
                section,open,class,extends,deriving,sorry,noncomputable},
  sensitive=true,
  morecomment=[l]{--},
  morecomment=[n]{/-}{-/},
  morestring=[b]",
}

% ── Lean listing style ───────────────────────────────────────────────
\lstdefinestyle{lean}{
  language=Lean4,
  backgroundcolor=\color{leanBg},
  basicstyle=\small\ttfamily,
  keywordstyle=\color{keyword}\bfseries,
  commentstyle=\color{comment}\itshape,
  breaklines=true,
  frame=single,
  framerule=0.4pt,
  xleftmargin=4pt,
  xrightmargin=4pt,
  aboveskip=6pt,
  belowskip=4pt,
  literate=
    {->}{{$\to$}}2
    {=>}{{$\Rightarrow$}}2
    {<|}{{$\langle$}}1
    {|>}{{$\rangle$}}1
    {==>}{{$\Longrightarrow$}}3
    { o. }{{$\;\circ\;$}}3
    {forall}{{$\forall$}}1
    {:=}{{$\,:=\,$}}2
    {!=}{{$\neq$}}2
    {<=}{{$\leq$}}2
    {>=}{{$\geq$}}2,
}

% ── Haskell listing style ───────────────────────────────────────────
\lstdefinestyle{haskell}{
  language=Haskell,
  backgroundcolor=\color{haskellBg},
  basicstyle=\small\ttfamily,
  keywordstyle=\color{keyword}\bfseries,
  commentstyle=\color{comment}\itshape,
  breaklines=true,
  frame=single,
  framerule=0.4pt,
  xleftmargin=4pt,
  xrightmargin=4pt,
  aboveskip=6pt,
  belowskip=4pt,
  literate=
    {->}{{$\to$}}2
    {=>}{{$\Rightarrow$}}2,
}

% ── Theorem environments ────────────────────────────────────────────
\theoremstyle{definition}
\newtheorem{defn}{Definition}[section]
\newtheorem{thm}[defn]{Theorem}
\newtheorem{prop}[defn]{Proposition}

\theoremstyle{remark}
\newtheorem*{fidelitynote}{Fidelity note}

% ── Convenience macros ───────────────────────────────────────────────
\newcommand{\lean}[1]{\texttt{#1}}      % inline Lean identifier
\newcommand{\hask}[1]{\texttt{#1}}      % inline Haskell identifier
\newcommand{\Ext}{\mathrm{Ext}}
\newcommand{\Pos}{\mathrm{Pos}}
\newcommand{\Shape}{\mathrm{Shape}}


\title{Human Verification: \texttt{Containers.Comonoid}\\[4pt]
  \large Directed Containers as Comonoids in $(\mathbf{Cont}, \triangleleft, I)$}
\author{Machine-generated from \texttt{lean/Comonoid.lean} for review by Neil Ghani}
\date{July 2026}

\begin{document}
\maketitle

\begin{abstract}
This document presents every definition and theorem \emph{statement} from the
Lean~4 module \lean{Containers.Comonoid} in three parallel forms: the original
Lean, a readable Haskell rendering, and precise mathematical English.  Proofs
are omitted --- they are machine-checked by Lean.  The purpose is to let a
human verifier confirm that the formalisation captures the intended mathematics.
\end{abstract}

\tableofcontents
\bigskip

%===================================================================
\section{Overview}
%===================================================================

The category $\mathbf{Cont}$ of containers carries a monoidal structure
$(\triangleleft, I)$, where the \emph{sequential composition} $C \triangleleft D$
is the container whose extension is $\mathbf{Ext}\,C \circ \mathbf{Ext}\,D$.  A
\textbf{comonoid} in $(\mathbf{Cont}, \triangleleft, I)$ is a container $C$
equipped with:
\begin{itemize}
  \item a \emph{counit morphism} $\varepsilon : C \to I$, and
  \item a \emph{comultiplication morphism} $\delta : C \to C \triangleleft C$,
\end{itemize}
satisfying the two counit laws and coassociativity, all stated as equalities of
container morphisms --- \emph{without ever passing to $[\mathbf{Set},\mathbf{Set}]$}.

The main theorem is that every directed container is canonically a
$\triangleleft$-comonoid, via \lean{DirectedContainer.toComonoid}.  The comonoid
laws correspond exactly to the directed-container axioms D1--D5:

\begin{center}
\begin{tabular}{lll}
  \textbf{Law} & \textbf{Comonoid} & \textbf{DC axioms} \\
  Left counit  & $\delta ; (\varepsilon \triangleleft C) = \lambda^{-1}$ &
    D1 (shapes) + D2 (positions) \\
  Right counit & $\delta ; (C \triangleleft \varepsilon) = \rho^{-1}$ &
    D3 \\
  Coassociativity & $\delta ; (\delta \triangleleft C) ; \alpha = \delta ; (C \triangleleft \delta)$ &
    D4 (shapes) + D5 (positions)
\end{tabular}
\end{center}

Everything is \textbf{Type}-level, Lean~4 core, no Mathlib.

%===================================================================
\section{Auxiliary lemmas}
%===================================================================

% --- 1. ContainerMorphism.ext_eq ----------------------------------
\subsection{Lemma: \lean{ContainerMorphism.ext\_eq}}

\begin{prop}[Pointwise extensionality for morphisms]
Two container morphisms $(f, f_P), (g, g_P) : C \to D$ are equal whenever:
\begin{itemize}
  \item their shape maps agree pointwise: $\forall s,\; f(s) = g(s)$, and
  \item their position maps agree after transporting along that pointwise
    equality: $\forall s, p,\; f_P(s, p) = g_P(s, (h_s)_* p)$,
    where $h_s : f(s) = g(s)$.
\end{itemize}
This is stated with a \emph{pointwise} shape equality deliberately, so that
the transport $h_s$ appears literally in the position obligation and can be
rewritten by \lean{Container.seq\_pos\_transport}.
\end{prop}

\paragraph{Lean.}
\begin{lstlisting}[style=lean]
theorem ContainerMorphism.ext_eq {C D : Container} {f g : C.Shape -> D.Shape}
    {fp : (s : C.Shape) -> D.Pos (f s) -> C.Pos s}
    {gp : (s : C.Shape) -> D.Pos (g s) -> C.Pos s}
    (hs : forall s, f s = g s) (hp : forall s p, fp s p = gp s (hs s cast p)) :
    (<|f, fp|> : ContainerMorphism C D) = <|g, gp|>
\end{lstlisting}

\paragraph{Haskell.}
\begin{lstlisting}[style=haskell]
-- Morphism extensionality: two morphisms with the same shape map (pointwise)
-- and compatible position maps are equal.  In Haskell this is just function
-- extensionality applied to both components.
\end{lstlisting}


% --- 2. Container.cast_symm_cast ----------------------------------
\subsection{Lemma: \lean{Container.cast\_symm\_cast}}

\begin{prop}[Round-trip transport cancellation]
For any type family $P : \alpha \to \mathbf{Type}$, equality $h : a = b$,
and element $x : P(a)$,
\[
  h^{-1}_*\bigl(h_* x\bigr) \;=\; x.
\]
Transporting along $h$ and back is the identity.
\end{prop}

\paragraph{Lean.}
\begin{lstlisting}[style=lean]
theorem Container.cast_symm_cast {alpha : Type} {P : alpha -> Type}
    {a b : alpha} (h : a = b) (x : P a) :
    (h.symm cast (h cast x : P b) : P a) = x
\end{lstlisting}

\paragraph{Haskell.}
\begin{lstlisting}[style=haskell]
-- The round-trip  coerce . coerce = id  holds by parametricity / type safety.
-- In Lean this requires an explicit proof because transport is not definitionally
-- involutive.
\end{lstlisting}


% --- 3. Container.seq_pos_transport --------------------------------
\subsection{Lemma: \lean{Container.seq\_pos\_transport}}

\begin{prop}[Transport on positions of sequential composition]
Let $C, E$ be containers, $t : E.\mathrm{Shape}$, and
$f, g : E.\mathrm{Pos}(t) \to C.\mathrm{Shape}$ inner-shape assignments with
$h : \forall u,\; f(u) = g(u)$.  Then transporting a position
$\langle u, q \rangle : (E \triangleleft C).\mathrm{Pos}\langle t, f \rangle$
along the induced equality acts componentwise:
\[
  (h_*)_*\langle u, q \rangle
  \;=\;
  \langle u,\; (h_u)_* q \rangle
  \;\in\;
  (E \triangleleft C).\mathrm{Pos}\langle t, g \rangle.
\]
That is, the outer position $u$ is unchanged, and only the inner position $q$
is transported along $h(u) : f(u) = g(u)$.
\end{prop}

\paragraph{Lean.}
\begin{lstlisting}[style=lean]
theorem Container.seq_pos_transport {C E : Container} {t : E.Shape}
    {f g : E.Pos t -> C.Shape} (h : forall u, f u = g u)
    (u : E.Pos t) (q : C.Pos (f u)) :
    (congrArg (Sigma.mk (beta := fun k => E.Pos k -> C.Shape) t) (funext h) cast
        (<|u, q|> : (E seq C).Pos <|t, f|>) : (E seq C).Pos <|t, g|>)
      = (<|u, (h u cast q : C.Pos (g u))|> : (E seq C).Pos <|t, g|>)
\end{lstlisting}

\paragraph{Haskell.}
\begin{lstlisting}[style=haskell]
-- The position type of (E seq C) over shape (t, f) is  (u : E.Pos t, C.Pos (f u)).
-- Transporting along a pointwise equality h : forall u. f u = g u
-- leaves u unchanged and coerces the inner position via h u.
-- No Haskell encoding is natural; this is purely a transport-bookkeeping lemma.
\end{lstlisting}

\begin{fidelitynote}
This lemma and \lean{ContainerMorphism.ext\_eq} are the two main tools for
managing dependent transports in this file.  They have no non-trivial Haskell
analogues, since Haskell lacks dependent types.  Their role is to isolate and
discharge all ``bureaucratic'' casts so that the mathematical content of the
comonoid laws is visible in the main proof of \lean{toComonoid}.
\end{fidelitynote}


%===================================================================
\section{Comonoid structure}
%===================================================================

% --- 4. Container.Comonoid ----------------------------------------
\subsection{Structure: \lean{Container.Comonoid}}

\begin{defn}[Comonoid in $(\mathbf{Cont}, \triangleleft, I)$]
A \textbf{comonoid} on a container $C$ consists of:
\begin{itemize}
  \item a counit morphism $\varepsilon : C \to I$,
  \item a comultiplication morphism $\delta : C \to C \triangleleft C$,
\end{itemize}
satisfying:
\begin{itemize}
  \item \textbf{Left counit:}
    $\delta \mathbin{;} (\varepsilon \triangleleft \mathrm{id}_C)
    = \lambda^{-1}_C$,
  \item \textbf{Right counit:}
    $\delta \mathbin{;} (\mathrm{id}_C \triangleleft \varepsilon)
    = \rho^{-1}_C$,
  \item \textbf{Coassociativity:}
    $\bigl(\delta \mathbin{;} (\delta \triangleleft \mathrm{id}_C)\bigr)
      \mathbin{;} \alpha_{C,C,C}
    = \delta \mathbin{;} (\mathrm{id}_C \triangleleft \delta)$.
\end{itemize}
Here $\lambda, \rho, \alpha$ are the left unitor, right unitor, and associator
of the monoidal category $(\mathbf{Cont}, \triangleleft, I)$, and $;$ is morphism
composition in $\mathbf{Cont}$.
\end{defn}

\paragraph{Lean.}
\begin{lstlisting}[style=lean]
structure Container.Comonoid (C : Container) where
  counit : ContainerMorphism C Container.I
  comult : ContainerMorphism C (C seq C)
  left_counit :
    comult.comp (Container.seq2 counit (ContainerMorphism.id C))
      = (Container.leftUnitor C).inv
  right_counit :
    comult.comp (Container.seq2 (ContainerMorphism.id C) counit)
      = (Container.rightUnitor C).inv
  coassoc :
    (comult.comp (Container.seq2 comult (ContainerMorphism.id C))).comp
        (Container.associator C C C).hom
      = comult.comp (Container.seq2 (ContainerMorphism.id C) comult)
\end{lstlisting}

\paragraph{Haskell.}
\begin{lstlisting}[style=haskell]
-- A comonoid in (Cont, seq, I):
-- class Comonoid c where
--   counit :: c -> I          -- shape: always (); pos: root
--   comult :: c -> c seq c     -- shape: (s, sub s); pos: shift s p q
-- Laws (stated as equations between natural transformations):
--   left_counit  : comult ; (counit seq id) = leftUnitor_inv
--   right_counit : comult ; (id seq counit) = rightUnitor_inv
--   coassoc      : (comult ; (comult seq id)) ; assoc = comult ; (id seq comult)
\end{lstlisting}

\begin{fidelitynote}
The laws are stated \emph{internally} in $\mathbf{Cont}$ --- as equalities of
container morphisms, not as equalities of natural transformations.  This is the
key distinction from the comonad statement in \lean{Containers.Directed}: there
the laws mention $\mathbf{Ext}$, here they do not.  Unfolding the definitions
shows the two formulations are equivalent, and the proof of \lean{toComonoid}
makes this explicit.
\end{fidelitynote}


%===================================================================
\section{The directed container is a comonoid}
%===================================================================

% --- 5. coCounit --------------------------------------------------
\subsection{Definition: \lean{DirectedContainer.coCounit}}

\begin{defn}[Counit morphism]
The \textbf{counit container morphism} $\varepsilon : C \to I$ of a directed
container $D$ sends every shape to $()$ and pulls back positions to the root:
\[
  \varepsilon.\mathrm{onShapes}(s) \;=\; (), \qquad
  \varepsilon.\mathrm{onPos}(s, \_) \;=\; o(s).
\]
\end{defn}

\paragraph{Lean.}
\begin{lstlisting}[style=lean]
def coCounit (D : DirectedContainer) : ContainerMorphism D.toContainer Container.I where
  onShapes := fun _ => ()
  onPos := fun s _ => D.root s
\end{lstlisting}

\paragraph{Haskell.}
\begin{lstlisting}[style=haskell]
coCounit :: DirectedContainer s p -> Morph s p () ()
coCounit d = Morph { onShapes = \_ -> (), onPos = \s _ -> root d s }
\end{lstlisting}


% --- 6. coComult --------------------------------------------------
\subsection{Definition: \lean{DirectedContainer.coComult}}

\begin{defn}[Comultiplication morphism]
The \textbf{comultiplication container morphism} $\delta : C \to C \triangleleft C$
of a directed container $D$ is:
\[
  \delta.\mathrm{onShapes}(s) \;=\; \langle s, s \mathop{\downarrow} - \rangle, \qquad
  \delta.\mathrm{onPos}(s, \langle p, q \rangle) \;=\; p \oplus q.
\]
The shape map sends $s$ to the pair $(s, s \downarrow -)$; the position map
combines a pair $\langle p, q \rangle$ of positions via shift.
\end{defn}

\paragraph{Lean.}
\begin{lstlisting}[style=lean]
def coComult (D : DirectedContainer) :
    ContainerMorphism D.toContainer (D.toContainer seq D.toContainer) where
  onShapes := fun s => <|s, D.sub s|>
  onPos := fun s pq => D.shift s pq.1 pq.2
\end{lstlisting}

\paragraph{Haskell.}
\begin{lstlisting}[style=haskell]
coComult :: DirectedContainer s p -> Morph s p (s, s -> p, p -> s) (p, p) p
-- Shape map:  s |-> (s, sub d s)
-- Pos map:    (s, (p, q)) |-> shift d s p q
\end{lstlisting}


% --- 7. toComonoid ------------------------------------------------
\subsection{Theorem: \lean{DirectedContainer.toComonoid}}

\begin{thm}[Every directed container is a $\triangleleft$-comonoid]
Let $D$ be a directed container.  Then the pair
$(\mathtt{coCounit}(D), \mathtt{coComult}(D))$ makes $D.\mathrm{toContainer}$
into a comonoid in $(\mathbf{Cont}, \triangleleft, I)$.  Specifically:
\begin{itemize}
  \item \textbf{Left counit}: follows from D1 (shapes) and D2 (positions), using
    \lean{Container.seq\_pos\_transport} to handle the dependent cast.
  \item \textbf{Right counit}: follows immediately from D3.
  \item \textbf{Coassociativity}: follows from D4 (shapes) and D5 (positions),
    using two applications of \lean{Container.seq\_pos\_transport} together with
    \lean{Container.cast\_symm\_cast} to cancel the D4 round-trip.
\end{itemize}
\end{thm}

\paragraph{Lean.}
\begin{lstlisting}[style=lean]
def toComonoid (D : DirectedContainer) : Container.Comonoid D.toContainer where
  counit := D.coCounit
  comult := D.coComult
  left_counit := ...   -- proof omitted
  right_counit := ...  -- proof omitted
  coassoc      := ...  -- proof omitted
\end{lstlisting}

\paragraph{Haskell.}
\begin{lstlisting}[style=haskell]
-- Every directed container (s, p, root, sub, shift) satisfying D1--D5
-- gives a comonoid in (Cont, seq, I):
--   counit = Morph (\_ -> ()) (\s _ -> root s)
--   comult = Morph (\s -> (s, sub s)) (\s (p,q) -> shift s p q)
-- and the comonoid laws follow from D1--D5.
\end{lstlisting}

\begin{fidelitynote}
The Lean definition is a \lean{def}, not a \lean{theorem}, because it constructs
a term of type \lean{Container.Comonoid D.toContainer}.  The proof fields
(\lean{left\_counit}, \lean{right\_counit}, \lean{coassoc}) are filled by tactic
blocks; these are omitted here.

The correspondence between comonoid laws and DC axioms is exact:
\[
\text{right counit} \leftrightarrow \text{D3}, \quad
\text{left counit} \leftrightarrow \text{D1}+\text{D2}, \quad
\text{coassociativity} \leftrightarrow \text{D4}+\text{D5}.
\]
The dependent casts arising from D1 and D4 (which are only propositional shape
equalities) are managed by the two auxiliary lemmas proved in
Section~\ref{sec:aux} above.
\end{fidelitynote}


%===================================================================
\section{Translation conventions}
%===================================================================

\begin{enumerate}
\item \textbf{Sequential composition $\triangleleft$.}  The operator
  \texttt{<|} (written \texttt{seq} in listings) in Lean is the sequential composition of containers.  Its
  extension satisfies
  $\mathbf{Ext}(C \triangleleft D) \cong \mathbf{Ext}\,C \circ \mathbf{Ext}\,D$.
  In Haskell it would be a type-level composition of functors.

\item \textbf{Unit container $I$.}  The unit $I$ has a single shape $()$
  and no positions.  Its extension is the identity functor.

\item \textbf{Monoidal structure.}  The left unitor $\lambda$, right unitor
  $\rho$, and associator $\alpha$ are container isomorphisms (pairs of inverse
  container morphisms) proved in \lean{Containers.Monoidal}.

\item \textbf{Laws stated internally.}  The comonoid laws are equations
  between \emph{container morphisms}, not between natural transformations.
  This is significant: the theorem ``directed containers are exactly comonoids
  in $(\mathbf{Cont}, \triangleleft, I)$'' holds at the internal level,
  sidestepping the functor category $[\mathbf{Set}, \mathbf{Set}]$ entirely.

\item \textbf{Proof bodies omitted.}  Only definitions and theorem
  \emph{statements} are translated.  Lean checks the proofs; the human
  verifier checks that the \emph{statements} match the intended mathematics.
\end{enumerate}

\end{document}
